\section{The Russian Turn}

Alexander Dugin has pointed out\footnote{\url{http://www.youtube.com/watch?v=quxWhp5y23E}} that Russia has had no Enlightenment. In the context of this survival \& ancient adversarial relationship towards the West, Nicholas Berdyaev\footnote{\url{http://www.berdyaev.com/berdiaev/berd\_lib/1913\_170.html}} is worth quoting in full (he is analyzing the Westernizer Trubetskoi):

\begin{quotationx}
In this world there is no mystical translucency. But after his devastating of the mystical meaning of all the earthly embodiments, Trubetskoy ought to arrive at a certain difficulty. \textbf{Rendered difficult for him is the question about the Church. In the Church he believes in nothing, and besides the Church he acknowledges nothing.} But the Church ought also to seem one of the ``earthly utopias'', that this is the ultimate recourse of that utopianism, which seeks the heavenly within the earthly, the Divine within the human, the other-worldly in the this-sidely. Trubetskoy ultimately to acknowledge, that upon earth, within earthly human history, that the Church also is an impossibility and an utopia, that it can appear only at the end of the world, only in the other-sidely. And as regards the Church he ought to affirm his idealistic thesis, i.e. to acknowledge it as being non-substantial, and as a norm, an ideal. I reiterate, that the critique of theocracy by Trubetskoy is excellent. But upon the path on which he stands, the critique of theocracy mustneeds also be extended to the Church, as upon an earthly utopian embodiment of the heavenly. He reduces the Church to the sacraments. But the Church has always striven to be more than the sacraments, it traverses the sphere of the earthly embodiment of the heavenly. The Catholic Church conceives of itself as a theocracy, the Kingdom of God upon earth. But the Orthodox Church also is not free of theocratic pretensions to order the world as regards itself. The separation of the Church from the state, which Trubetskoy desires, is also a diminution within the Church of earthly embodiments of the heavenly, a diluted churchly utopianism. Must needs there not be upon this path a refusal of embodiment of the heavenly into the earthly, until the very end of the world? Trubetskoy subsequently defends the secularisation of the whole of life and all of culture, and this mustneeds be acknowledged as very powerful on his side. It is necessary to renounce the lie and the sham of a Christian state, a Christian science, a Christian culture, etc, etc. The process of secularisation has an inner significance. But the whole of life and culture ought anew to become religious. Trubetskoy does not uncover the path to this, he does not point out the means for the spiritualisation of life. The Church also is the path of the making Divine the world. But in the metaphysics of Trubetskoy there is no place for the existence of the Church upon the earth. The Church as it were ought to remain in heaven.

He demolishes all the earthly utopias, but together with this he demolishes also the Church, as resulting from an utopia. Trubetskoy points out the whole difficulty of the existence of a Christian dominion. There is only one Christian dominion --- this is the domain of the birds of the air and the lilies of the fields. Freedom from care is the Gospel testament. But upon earth it is impossible to be free from care, it is impossible to live, as do the lilies of the field and the birds of the air, upon the earth it is necessary to be prudent of domain. How to get out of this? It is indeed impossible to put off resolution of the contradictions of life until the end of the world. Do they not take too lightly the burden of the religious antinomy of life? It is necessary to accept life to the end, sacrificially and tragically. There is already therein no justification for the religiously neutral, for an external to God sociability and an external to God culture. Trubetskoy was more in the right than Bulgakov, when he rejects dominion from a Christian point of view and does not reckon it a task Divine. Bulgakov indeed has transferred over upon heaven his own economics, his own sweat. But the truth of Bulgakov is in this, that he sees the tormentiveness of the religious problem of economics and he does not reckon it possible to remain on the soil of religious neutrality. It mustneeds be recognised, that religiously there is nothing that is neutral, there is nothing that is outside religion.
\end{quotationx}


What is interesting here is the critique of the Church -- the Church becomes (if it accepts the separation of Church \& State) nothing more than a very fancy and rarefied utopia. Since Nature is not Divine, nor is there such nonsense as a World-Soul, and since all else of interest occurs within the compartments of secular work (by definition, whatever is un-Church), then ``the Church'' is nothing but a holding action, a word we place for something we do not have. Or, more pointedly, a Utopia. This (of course) is the default political position of conservatism today (something which Evola notes in the essay on Junger, \emph{The Gordian Knot}\footnote{\url{http://www.centrostudilaruna.it/the-gordian-knot.html}}).

\begin{quotex}
Now it is important to point out that wherever forces belonging to the first of the three domains emerge and break forth, only the possibilities of the third domain can really face them. Any attempt to stem on the basis of forces and values of the intermediate zone can only be precarious, provisional and relative.
\end{quotex}


As the ``Church'' of what ``used to be'', conservatism merely attempts to slow or thwart the move toward the abyss. Men like Evola \& Steiner attempted to erect programs which would remedy this fault, but nothing has (as of yet) materialized (the question has been raised at Gornahoor, of a lack of courage).

The systematizers failed, but so did the prophets and statesmen, like Russell Kirk. Western fighters in this intellectual strife may well look outside of the West to find the crack in the mirror.

If we go back to Berdyaev, he here puts his finger (with Russian characteristic flair) on the problem with the Enlightenment ``mindset''. Seeing the contrast between secularism \& theocracy, the Russian soul either prefers to risk theocracy (``Third Rome'' and ``Slavophilism'') or to endeavor a Third Testament (The Church of John). This is why Gornahoor has spoken of a revival or preservation of hermeticism \& counter-revolution within Russia. Of course, many Russians attempted to secularize, but the attempt upon the soul (in general) failed.

A friend of mine (Tor, from Dunedain.net) writes

\begin{quotex}
\textbf{This is why we turn again to Holy Russia, because the madmen there have never stopped weeping and, if they have not picked up the sword, it is because their tear-rivers will serve to forge new ones.} There is a death beyond dying that can be approached in the living that is beyond living and this mode, if abandoned in the new age, is approachable by all who have that little mite that can move mountains. Yes, it is faith that I speak of, and this mite is more powerful than all armies arrayed against it, all manner of worldly men that have some stock in old castles and feudal titles or yachts and mighty steel towers, and who scoff at the world and at man and perceive themselves to be gods. Yeah, verily, this is the spirit not of those who were simply crushed by the church for their herbal potions, but of the martyrs who willingly marched under a blood-red banner and gave their lives, all and total, for the new life that would be theirs whither and whether they fell.
\end{quotex}


It may be that within the belly of the beast (whether Soviet Russia or post-capitalist America) is found the strongest chance to preserve one's self \& to (therefore) resist without resisting, that which seems irresistible.

In this struggle, every true Christian will have to be transmuted to a ``hermetic'', \& every esotericist will have to strive to attain the martyr spirit. When the two have become one, the lightning shall come from the East to the West.


\flrightit{Posted on 2011-07-23 by Logres}

